% Generated by tools/edn_latex.py from reports/edn.md.
% Do not edit: change reports/edn.md and run `make edn`.
\ednentry{
  id = {EDN-21},
  title = {Report the upstream licence contradiction instead of resolving it},
  date = {2026-09-29},
  milestone = {M1: Project Inception},
  topic = {Data Licensing and Provenance},
  participants = {@lukas2510},
  decision = {The dataset card reports both licence statements that the Zenodo record makes about the AutoScout24 dataset, and names the contradiction, rather than picking one or asking the author to resolve it. The YAML frontmatter carries \texttt{license: mit} because that is the record's structured field, and the body explains the discrepancy.},
  rationale = {\emph{Alternatives considered:}
\begin{itemize}[leftmargin=*, topsep=2pt]
\item \textbf{Option A (chosen): report both statements and flag the contradiction.}
\newline Pros: it is the only accurate description of what the source actually says; it costs nothing, because both readings permit non-commercial academic use; it warns anyone who later reuses our work for a commercial purpose, which is exactly where the difference would start to matter; it demonstrates the provenance care the course rewards.
\newline Cons: a reader has to absorb an ambiguity rather than a single answer; the frontmatter still has to commit to one identifier, so the card is slightly inconsistent with itself unless the body is read.
\item \textbf{Option B: report MIT only.}
\newline Pros: simplest; matches the machine-readable metadata field that tooling and aggregators would read; MIT is the more permissive of the two, so it cannot under-claim our own rights.
\newline Cons: silently discards a restriction the author wrote on the same record; would mislead a future reader who wanted to use the data commercially; asserts a certainty we do not have.
\item \textbf{Option C: report the restrictive prose only.}
\newline Pros: the conservative reading, so it cannot over-claim rights.
\newline Cons: contradicts the structured licence field; ``research, educational or analytical'' is not a recognised licence, so it is not expressible in the card's frontmatter; over-restricts what may genuinely be MIT.
\item \textbf{Option D: ask the author to resolve it before finishing the card.}
\newline Pros: would produce a definitive answer and a citable clarification.
\newline Cons: blocks a sprint deliverable on a stranger's response time with no deadline; disproportionate for a university project whose use is permitted under either reading.
\end{itemize}
\par
\emph{Why:} The contradiction is in the source, not in our reading of it: the Zenodo record's structured licence field says MIT while the author's description on the same record says ``You are welcome to use this dataset for research, educational, or analytical purposes''. Since our use is non-commercial academic work, both readings permit it, so resolving the ambiguity would change nothing about what we are allowed to do. Documenting it honestly is therefore strictly better than choosing a side, and materially better than blocking the deliverable on an email. Option D was explicitly declined by the team as disproportionate for the size of the project.},
  ai = {Information seeking, Alternative generation, Alternative assessment, Recommendation},
  response = {Accepted with modifications},
  assessment = {AI found the contradiction while re-validating PR \#22, and in doing so corrected its own earlier review, which had asserted that the card's restrictive wording was simply wrong and contradicted the MIT licence recorded in the project brief. Fetching the Zenodo API record showed that both statements come from the same source, so the card's author had not made a mistake and the project brief's flat ``License: MIT'' was itself incomplete. AI laid out the four options above and recommended option A. The team accepted the recommendation but rejected the accompanying suggestion to contact the author, on the grounds that it is disproportionate for a university project.},
  aievidence = {Claude Code session on 2026-09-29 while reviewing PR \#22. AI fetched \texttt{https:/\allowbreak{}/\allowbreak{}zenodo.\allowbreak{}org/\allowbreak{}api/\allowbreak{}records/\allowbreak{}17643343} and reported that \texttt{metadata.\allowbreak{}license.\allowbreak{}id} is \texttt{mit-\allowbreak{}license} while the description text welcomes use for research, education and analysis, which reads as narrower than MIT; it then flagged the choice as EDN-worthy rather than deciding it. Lukas replied, in German, that the team should not be pedantic and should simply take the recommendation, since this is a small university-course project, which settled option A and ruled out option D.},
  evidence = {\href{https://github.com/mlops-2627q1-mds-upc/MLOPS_Recommenditos/pull/22}{PR \#22}, the Licensing section of \href{https://github.com/mlops-2627q1-mds-upc/MLOPS_Recommenditos/blob/main/docs/docs/dataset-card.md}{the dataset card}, Zenodo record \href{https://doi.org/10.5281/zenodo.17643343}{10.5281/zenodo.17643343}.},
}
